\section{Proofs for Structural Properties of MEC}
\label{appendix:structural-properties-proofs}

\subsection{Convexity}
\subsubsection{Proof of Theorem~\ref{convexity}}
\label{appendix:proof-convexity}
\convexityThm*
\begin{proof}
From Def.~\ref{mec:def}, \(\beta \geq 1\). For any fixed \(x\in \support(\pi)\), consider 
\[
g_x(y)=|x-y|^\beta.
\]
It is a well-known result (see \cite{Rockafellar1970}) that \(g_x(y)\) is convex on \(\Reals\) for \(\beta\ge1\) and strictly convex for \(\beta>1\). Both convexity and strict convexity are preserved under multiplication by strictly positive constants. Hence, since for \(x\in \support(\pi)\) we have \(\pi(x)^\alpha>0\), the function 
\[
f_x(y)=\pi(x)^\alpha\, |x-y|^\beta
\]
is convex and, if \(\beta>1\), strictly convex as well. As \(\EC(y)\) is a finite sum of the functions \(f_x(y)\), and (strict) convexity is preserved under finite sums, it follows that \(\EC\) is convex for \(\beta\ge1\) and strictly convex for \(\beta>1\).
\end{proof}

\subsection{Characterization and Special Cases for MEC}
\subsubsection{Proof of Theorem~\ref{thm:mec-Lbeta}}
\label{appendix:proof-mec-lbeta}
\mecLbeta*
\begin{proof}
For $\alpha\ge 1$ and $\beta\ge 1$, the effort function is
\[
\EC(y)=\sum_{x\in\support(\pi)} \pi(x)^\alpha |x-y|^\beta
=\sum_{x\in\support(\pi)} \pi^\alpha(x)\,|x-y|^\beta.
\]
Factoring out $\Mass{\pi^\alpha}$ yields
\[
\EC(y)=\Mass{\pi^\alpha}\sum_{x\in\support(\pi)} \overline{\pi^\alpha}(x)\,|x-y|^\beta.
\]
Taking minima over $y\in[0,1]$ gives
\[
\MEC[\alpha,\beta](\pi)
=\min_{y\in[0,1]}\EC(y)
=\Mass{\pi^\alpha}\min_{y\in[0,1]}\sum_{x\in\support(\pi)} \overline{\pi^\alpha}(x)\,|x-y|^\beta
=\Mass{\pi^\alpha}\cdot D_\beta(\overline{\pi^\alpha}),
\]
as claimed.
\end{proof}

\subsubsection{Proof of Corollary~\ref{cor:mec-special-cases}}
\label{appendix:proof-cor-mec-special-cases}
\specialCases*
\begin{proof}
Let $\rho=\overline{\pi^\alpha}$. If $\beta=1$, the function $y\mapsto \sum_x \rho(x)\,|x-y|$ is minimized at any
median $m$ of $\rho$, hence
\[
\min_y \sum_x \rho(x)\,|x-y|=\sum_x \rho(x)\,|x-m|
=\mathrm{MAD}_{\mathrm{m}}(\rho).
\]
If $\beta=2$, the function $y\mapsto \sum_x \rho(x)\,(x-y)^2$ is minimized at
the mean $\mu$ of $\rho$, and the minimum equals
\[
\min_y \sum_x \rho(x)\,(x-y)^2=\sum_x \rho(x)\,(x-\mu)^2=\mathrm{Var}(\rho).
\]
Multiplying by $\Mass{\pi^\alpha}$ yields the stated identities.
\end{proof}

\subsection{Symmetry}

\subsubsection{Proof of Lemma~\ref{lem:mec-reflection}}
\label{appendix:proof-mec-reflection}

\mecReflection*
\begin{proof}
Let
\[
I_c \defsymbol [0,1]\cap[2c-1,2c].
\]
By definition of reflection,
\[
\support(\pi^{\mathrm{ref}(c)})
=
\{\,2c-x \mid x\in\support(\pi)\,\}
\]
and
\[
\pi^{\mathrm{ref}(c)}(2c-x)\defsymbol \pi(x)
\qquad\text{for all }x\in\support(\pi).
\]
Since $c$ is an admissible reflection center for $\pi$, for every
$x\in\support(\pi)$ we have $2c-x\in[0,1]$, hence
\[
2c-1\le x\le 2c.
\]
Therefore
\[
\support(\pi)\subseteq I_c
\qquad\text{and}\qquad
\support(\pi^{\mathrm{ref}(c)})\subseteq I_c.
\]
For every $y\in I_c$ we have $2c-y\in I_c\subseteq[0,1]$, and
\[
\begin{aligned}
\EC[\alpha,\beta][\pi^{\mathrm{ref}(c)}](y)
&=
\sum_{x\in\support(\pi^{\mathrm{ref}(c)})}
\pi^{\mathrm{ref}(c)}(x)^\alpha |x-y|^\beta \\
&=
\sum_{z\in\support(\pi)}
\pi(z)^\alpha |(2c-z)-y|^\beta \\
&=
\sum_{z\in\support(\pi)}
\pi(z)^\alpha |z-(2c-y)|^\beta \\
&=
\EC(2c-y).
\end{aligned}
\]
Since $\support(\pi)\subseteq I_c$, the minimum of $\EC$ over $[0,1]$
is attained in $I_c$, and the same holds for $\EC[\alpha,\beta][\pi^{\mathrm{ref}(c)}]$.
Because $y\mapsto 2c-y$ is a bijection from $I_c$ onto itself, we obtain
\[
\MEC[\alpha,\beta](\pi^{\mathrm{ref}(c)})
=
\MEC[\alpha,\beta](\pi).
\]
\end{proof}

\subsubsection{Proof of Corollary~\ref{cor:mec-shift}}
\label{appendix:proof-mec-shift}

\mecShift*
\begin{proof}
Observe that a shift can be expressed as the composition of two reflections.
If $c$ and $c+\tfrac{t}{2}$ are admissible reflection centers for the relevant
distributions, then
\[
\pi^{\mathrm{shift}(t)}
=
\left(\pi^{\mathrm{ref}(c)}\right)^{\mathrm{ref}(c+t/2)},
\]
since for every $x\in\support(\pi)$,
\[
2\!\left(c+\frac{t}{2}\right)-(2c-x)=x+t.
\]
In particular, if
\[
c=\frac{x_1+x_n}{2},
\]
where $x_1=\min\support(\pi)$ and $x_n=\max\support(\pi)$, then $c$ is an
admissible reflection center for $\pi$; and if $t$ is an admissible shift,
then $c+\tfrac{t}{2}$ is an admissible reflection center for
$\pi^{\mathrm{ref}(c)}$.
\end{proof}

\subsubsection{Proof of Lemma~\ref{lem:symmetry-mec}}
\label{appendix:proof-symmetry-mec}

\symmetryMec*
\begin{proof}
Recall that
\[
\EC(y)=\sum_{x\in\support(\pi)} \pi(x)^\alpha |x-y|^\beta.
\]
Since $\pi$ is symmetric about $c$, for every $x\in\support(\pi)$ we have
$2c-x\in\support(\pi)$ and
\[
\pi(x)=\pi(2c-x).
\]
Hence every point of $\support(\pi)$ distinct from $c$ belongs to a symmetric
pair $\{x,2c-x\}$, and if $c\in\support(\pi)$ it forms a singleton.

We first show that $\EC$ is symmetric about $c$. Let $t$ be such that
$c\pm t\in[0,1]$. For each symmetric pair $\{x,2c-x\}$, using
$\pi(x)=\pi(2c-x)$, we have
\[
\begin{aligned}
&\pi(x)^\alpha |x-(c+t)|^\beta
+\pi(2c-x)^\alpha |(2c-x)-(c+t)|^\beta \\
&\qquad
=\pi(x)^\alpha |x-(c+t)|^\beta
+\pi(x)^\alpha |x-(c-t)|^\beta .
\end{aligned}
\]
Interchanging $t$ and $-t$ leaves this quantity unchanged, so the
contribution of each symmetric pair to $\EC(c+t)$ equals its contribution
to $\EC(c-t)$. If $c\in\support(\pi)$, then the singleton contributes
\[
\pi(c)^\alpha |c-(c+t)|^\beta
=
\pi(c)^\alpha |t|^\beta
=
\pi(c)^\alpha |c-(c-t)|^\beta,
\]
which is also the same at $c+t$ and $c-t$. Summing over all support points,
we obtain
\[
\EC(c+t)=\EC(c-t).
\]
Therefore $\EC$ is symmetric about $c$.

If $\beta>1$, then by Theorem~\ref{convexity}, $\EC$ is strictly convex on
$[0,1]$, so it has a unique optimal consensus point. Since $\EC$ is symmetric about $c$,
that unique optimal consensus point must be $c$. Hence
\[
\MEC[\alpha,\beta](\pi)=\EC(c).
\]

Assume now that $\beta=1$. By Corollary~\ref{cor:mec-special-cases}, the
optimal consensus points of $\pi$ are precisely the medians of $\overline{\pi^\alpha}$. Since
\[
\pi(x)^\alpha=\pi(2c-x)^\alpha
\qquad\text{for all }x\in[0,1],
\]
the distribution $\pi^\alpha$ is also symmetric about $c$. Therefore the
total $\pi^\alpha$-mass strictly to the left of $c$ equals the total
$\pi^\alpha$-mass strictly to the right of $c$, and thus both are at most
$\Mass{\pi^\alpha}/2$. Hence $c$ is a median of $\pi^\alpha$, so $\EC(c)$
attains the minimum. Therefore again
\[
\MEC[\alpha,\beta](\pi)=\EC(c).
\]

Thus, in all cases,
\[
\MEC[\alpha,\beta](\pi)
=
\EC(c)
=
\sum_{x\in\support(\pi)} \pi(x)^\alpha |x-c|^\beta.
\]
\end{proof}

\subsubsection{Proof of Theorem~\ref{th:mec-pi-star}}
\label{appendix:proof-mec-pi-star}

\mecPiStar*
\begin{proof}
Since $\pi^*(0)=\pi^*(1)=\Mass{\pi^*}/2$, the distribution $\pi^*$ is symmetric about $1/2$. Hence, by Lemma~\ref{lem:symmetry-mec},
\[
\MEC(\pi^*)=\EC[\alpha,\beta][\pi^*](1/2).
\]
Evaluating the effort at $1/2$ gives
\[
\EC[\alpha,\beta][\pi^*](1/2)
=
\left(\frac{\Mass{\pi^*}}{2}\right)^\alpha\!\left(\frac12\right)^\beta
+
\left(\frac{\Mass{\pi^*}}{2}\right)^\alpha\!\left(\frac12\right)^\beta
=
2\left(\frac{\Mass{\pi^*}}{2}\right)^\alpha\left(\frac12\right)^\beta.
\]
\end{proof}

\subsection{Population Scaling}

\subsubsection{Proof of Theorem~\ref{th:multbyscalar}}
\label{appendix:proof-population-scaling}

\multByScalar*
\begin{proof}
By definition,
\[
\MEC(\lambda\pi)
=
\min_{y\in[0,1]}
\sum_{i=1}^{n}(\lambda\pi_i)^{\alpha}|x_i-y|^{\beta}.
\]
Since
\[
(\lambda\pi_i)^{\alpha}=\lambda^{\alpha}\pi_i^{\alpha},
\]
it follows that
\[
\MEC(\lambda\pi)
=
\min_{y\in[0,1]}
\lambda^{\alpha}\sum_{i=1}^{n}\pi_i^{\alpha}|x_i-y|^{\beta}.
\]
Because $\lambda^{\alpha}$ does not depend on $y$, it can be factored out of the minimum:
\[
\MEC(\lambda\pi)
=
\lambda^{\alpha}
\min_{y\in[0,1]}
\sum_{i=1}^{n}\pi_i^{\alpha}|x_i-y|^{\beta}.
\]
Therefore,
\[
\MEC(\lambda\pi)=\lambda^{\alpha}\MEC(\pi).
\]
\end{proof}

\subsubsection{Proof of Corollary~\ref{cor:condition-H-normalization}}
\label{appendix:proof-condition-H-normalization}

\conditionH*
\begin{proof}
(1) By definition of normalization,
\[
\overline{\pi}=\frac{1}{\Mass{\pi}}\,\pi.
\]
Applying Th.~\ref{th:multbyscalar} with
\[
\lambda=\frac{1}{\Mass{\pi}},
\]
we obtain
\[
\MEC(\overline{\pi})
=
\MEC\!\left(\frac{1}{\Mass{\pi}}\pi\right)
=
\left(\frac{1}{\Mass{\pi}}\right)^\alpha \MEC(\pi).
\]
Therefore,
\[
\MEC(\overline{\pi})
=
\frac{\MEC(\pi)}{(\Mass{\pi})^\alpha},
\]
as claimed.

(2) From Th.~\ref{th:multbyscalar} we have
\[
\MEC(\lambda\pi)=\lambda^{\alpha}\MEC(\pi)
\qquad\text{and}\qquad
\MEC(\lambda\pi')=\lambda^{\alpha}\MEC(\pi').
\]
Since $\lambda^{\alpha}>0$, multiplying the inequality
\[
\MEC(\pi)>\MEC(\pi')
\]
by $\lambda^{\alpha}$ preserves the order, yielding
\[
\MEC(\lambda\pi)>\MEC(\lambda\pi').
\]
Hence MEC satisfies Condition~H.
\end{proof}


